\documentclass[aspectratio=169,10pt]{beamer}
% ---------------------------------------------------------------
% R3.07 SQL dans un langage de programmation -- preambule commun des diapositives
% Charte derivee de R5.A.08 Qualite de developpement (metropolis, dark tech)
% Palette IUT partagee (violet du logo + accent or)
% Utilisation :
%   \documentclass[aspectratio=169,10pt]{beamer}
%   % ---------------------------------------------------------------
% R3.07 SQL dans un langage de programmation -- preambule commun des diapositives
% Charte derivee de R5.A.08 Qualite de developpement (metropolis, dark tech)
% Palette IUT partagee (violet du logo + accent or)
% Utilisation :
%   \documentclass[aspectratio=169,10pt]{beamer}
%   % ---------------------------------------------------------------
% R3.07 SQL dans un langage de programmation -- preambule commun des diapositives
% Charte derivee de R5.A.08 Qualite de developpement (metropolis, dark tech)
% Palette IUT partagee (violet du logo + accent or)
% Utilisation :
%   \documentclass[aspectratio=169,10pt]{beamer}
%   \input{../../latex/common/r307-slides-preamble.tex}
%   \graphicspath{{../../latex/common/assets/}{img/}{figs/}}
% ---------------------------------------------------------------
\usetheme{metropolis}
\usepackage{appendixnumberbeamer}

\usepackage[T1]{fontenc}
\usepackage[utf8]{inputenc}
\usepackage[french]{babel}
\usepackage[scale=0.94]{tgheros}
\renewcommand{\familydefault}{\sfdefault}
\usepackage{amsmath,amssymb,amsfonts}
\usepackage{graphicx}
\usepackage{booktabs}
\usepackage{multirow}
\usepackage{tikz}
\usetikzlibrary{shapes.geometric, shapes, arrows, arrows.meta, positioning, calc, fit}
\usepackage{xcolor}
\usepackage{listings}
\usepackage[skins,listings]{tcolorbox}
\usepackage{qrcode}
\usepackage{hyperref}

% ---------------------------------------------------------------
% PALETTE (dark-mode tech: near-black canvas, violet (logo IUT) + gold accents)
% ---------------------------------------------------------------
\definecolor{qabg}{HTML}{0E1016}      % canvas background
\definecolor{qapanel}{HTML}{181B24}   % card / block panel
\definecolor{qapanelalt}{HTML}{20242F}% secondary panel tint
\definecolor{qaline}{HTML}{2B303D}    % hairline rules / table rules
\definecolor{qatext}{HTML}{ECEEF2}    % primary text (off-white)
\definecolor{qamuted}{HTML}{8B93A6}   % secondary / caption text
\definecolor{qaaccent}{HTML}{D93ED7}  % primary accent -- violet du logo IUT, eclairci pour fond sombre
\definecolor{qaaccent2}{HTML}{F2B84B} % secondary accent -- or/ambre
\definecolor{qacodekw}{HTML}{C792EA}  % code : mots-cles
\definecolor{qacodestr}{HTML}{C3E88D} % code : chaines
\definecolor{qacodecom}{HTML}{7F8AA3} % code : commentaires
\definecolor{qacodeann}{HTML}{F2B84B} % code : annotations / javadoc

\setbeamercolor{background canvas}{bg=qabg}
\setbeamercolor{normal text}{fg=qatext}
\setbeamercolor{palette primary}{bg=qabg, fg=qamuted}
\setbeamercolor{frametitle}{bg=qabg, fg=qatext}
\setbeamercolor{progress bar}{fg=qaaccent, bg=qaline}
\setbeamercolor{alerted text}{fg=qaaccent2}
\setbeamercolor{block title}{bg=qapanel, fg=qaaccent}
\setbeamercolor{block body}{bg=qapanel, fg=qatext}
\setbeamercolor{block title alerted}{bg=qapanel, fg=qaaccent2}
\setbeamercolor{block body alerted}{bg=qapanel, fg=qatext}
\setbeamercolor{block title example}{bg=qapanelalt, fg=qaaccent}
\setbeamercolor{block body example}{bg=qapanelalt, fg=qatext}
\setbeamercolor{example text}{fg=qaaccent}
\setbeamercolor{itemize item}{fg=qaaccent}
\setbeamercolor{itemize subitem}{fg=qaaccent2}
\setbeamercolor{enumerate item}{fg=qaaccent}
\setbeamercolor{enumerate subitem}{fg=qaaccent2}
\setbeamercolor{caption}{fg=qamuted}
\setbeamercolor{page number in head/foot}{fg=qamuted, bg=qabg}
\setbeamercolor{title}{fg=qatext}
\setbeamercolor{subtitle}{fg=qamuted}
\setbeamercolor{author}{fg=qatext}
\setbeamercolor{institute}{fg=qamuted}
\setbeamercolor{date}{fg=qamuted}
\setbeamercolor{section in toc}{fg=qatext}
\setbeamercolor{subsection in toc}{fg=qamuted}

\setbeamertemplate{blocks}[rounded][shadow=false]
\setbeamertemplate{itemize item}{\raisebox{0.15em}{\scriptsize\textcolor{qaaccent}{\ensuremath{\blacktriangleright}}}}
\setbeamertemplate{itemize subitem}{\raisebox{0.1em}{\tiny\textcolor{qaaccent2}{\ensuremath{\bullet}}}}

\setbeamerfont{frametitle}{series=\bfseries, size=\Large}
\setbeamerfont{title}{series=\bfseries, size=\Large}
\setbeamerfont{subtitle}{size=\normalsize}
\setbeamerfont{institute}{size=\scriptsize}
\setbeamerfont{date}{size=\scriptsize}
\setbeamerfont{author}{size=\small}

\metroset{sectionpage=none, subsectionpage=none, numbering=fraction, progressbar=frametitle}
\setbeamertemplate{section in toc}[sections numbered]

% ---------------------------------------------------------------
% LISTINGS -- theme sombre (lstlisting reste utilisable tel quel)
% \lstcode : police + couleur de base, a reutiliser dans les
%            surcharges locales (ex. basicstyle=\lstcode\scriptsize)
% ---------------------------------------------------------------
\newcommand{\lstcode}{\ttfamily\color{qatext}}

\lstdefinelanguage{TypeScript}{
  keywords={typeof, new, true, false, catch, function, return, null, switch, var,
    if, in, while, do, else, case, break, const, let, interface, class, extends,
    implements, enum, export, import, this, super, constructor, public, private,
    protected, static, async, await, throw, throws, try, finally, void},
  ndkeywords={boolean, any, number, string, int, long, double},
  sensitive=true,
  comment=[l]{//},
  morecomment=[s]{/*}{*/},
  morestring=[b]',
  morestring=[b]"
}

\lstset{
  language=Java,
  basicstyle=\lstcode\footnotesize,
  keywordstyle=\color{qacodekw}\bfseries,
  ndkeywordstyle=\color{qaaccent2}\bfseries,
  identifierstyle=\color{qatext},
  stringstyle=\color{qacodestr},
  commentstyle=\color{qacodecom}\itshape,
  morecomment=[s][\color{qacodeann}]{/**}{*/},
  backgroundcolor=\color{qapanelalt},
  rulecolor=\color{qaline},
  frame=single,
  framesep=3pt,
  xleftmargin=4pt,
  xrightmargin=4pt,
  keepspaces=true,
  columns=fullflexible,
  showstringspaces=false,
  breaklines=true,
  tabsize=2,
  escapeinside={(*@}{@*)},
}

% ---------------------------------------------------------------
% CARD-STYLE BOXES
% ---------------------------------------------------------------
\tcbset{
  figcard/.style={colback=white, colframe=white, boxrule=0pt, arc=4pt,
    left=3pt, right=3pt, top=3pt, bottom=3pt, boxsep=0pt,
    halign=center, valign=center},
  datacard/.style={colback=qapanel, colframe=qaline, boxrule=0.4pt, arc=3pt,
    left=8pt, right=8pt, top=6pt, bottom=6pt, coltext=qatext,
    before upper={\color{qatext}}, halign=center},
  codebox/.style={colback=qapanelalt, colframe=qaline, boxrule=0.4pt, arc=3pt,
    left=8pt, right=8pt, top=5pt, bottom=5pt, listing only,
    listing options={basicstyle=\ttfamily\small\color{qaaccent},
      breaklines=true, showstringspaces=false, columns=fullflexible,
      keepspaces=true, frame=none, backgroundcolor={}}}
}
\newtcolorbox{figcard}[1][]{figcard,#1}
\newtcolorbox{datacard}{datacard}
\newtcblisting{codebox}{codebox}

% ---------------------------------------------------------------
% HANDY MACROS
% ---------------------------------------------------------------
\newcommand{\hl}[1]{\alert{#1}}
\newcommand{\source}[1]{{\scriptsize\color{qamuted}#1}}

% Sommaire automatique en debut de section
\AtBeginSection[]{%
  \begin{frame}<beamer>{Plan du cours}
    \tableofcontents[currentsection]
  \end{frame}}

% Page de titre : \rtitleslide{numero de chapitre}{titre du chapitre}
% Les auteurs sont repris de \rauthors (modifiable avant \begin{document})
\newcommand{\rauthors}{Selena Lamari et Chouki Tibermacine}
\newcommand{\rtitleslide}[2]{%
\begin{frame}[plain,noframenumbering]
  \begin{center}
    \vspace{2mm}
    {\color{qaaccent}\rule{2.1cm}{1.2pt}}\\[6mm]
    {\Large\bfseries\color{qatext} R3.07}\\[1mm]
    {\Large\bfseries\color{qatext} SQL dans un langage de programmation}\\[5mm]
    {\large\color{qamuted} Chapitre #1}\\[1mm]
    {\large\color{qaaccent} #2}\\[6mm]
    \includegraphics[height=2.1cm]{iut_logo_transparent.png}

    \vspace{4mm}

    {\small\color{qatext} \rauthors}\\[1mm]
    {\scriptsize\color{qamuted} BUT 2 Informatique --- IUT de Vannes --- Universit\'e Bretagne Sud}
  \end{center}
\end{frame}
}

% ---- Diapo de fin « Questions ? » (même modèle que R5.A.08) ----
\newcommand{\rcontact}{Chouki.Tibermacine@univ-ubs.fr}
\newcommand{\rquestionslide}{%
\begin{frame}{Questions ?}
  \begin{center}
    \vspace{2mm}
    \begin{tikzpicture}
      \node[regular polygon, regular polygon sides=6, draw=qaline, thick,
        minimum size=4.6cm, rotate=90] (hexin) {};
      \node[regular polygon, regular polygon sides=6, draw=qaaccent, dashed,
        thick, minimum size=6.2cm, rotate=90] (hexout) {};
      \node[font=\fontsize{92}{92}\selectfont\bfseries, color=qaaccent2] at (hexin) {?};
    \end{tikzpicture}\\[5mm]
    {\small\color{qamuted} \rcontact}
  \end{center}
\end{frame}}

%   \graphicspath{{../../latex/common/assets/}{img/}{figs/}}
% ---------------------------------------------------------------
\usetheme{metropolis}
\usepackage{appendixnumberbeamer}

\usepackage[T1]{fontenc}
\usepackage[utf8]{inputenc}
\usepackage[french]{babel}
\usepackage[scale=0.94]{tgheros}
\renewcommand{\familydefault}{\sfdefault}
\usepackage{amsmath,amssymb,amsfonts}
\usepackage{graphicx}
\usepackage{booktabs}
\usepackage{multirow}
\usepackage{tikz}
\usetikzlibrary{shapes.geometric, shapes, arrows, arrows.meta, positioning, calc, fit}
\usepackage{xcolor}
\usepackage{listings}
\usepackage[skins,listings]{tcolorbox}
\usepackage{qrcode}
\usepackage{hyperref}

% ---------------------------------------------------------------
% PALETTE (dark-mode tech: near-black canvas, violet (logo IUT) + gold accents)
% ---------------------------------------------------------------
\definecolor{qabg}{HTML}{0E1016}      % canvas background
\definecolor{qapanel}{HTML}{181B24}   % card / block panel
\definecolor{qapanelalt}{HTML}{20242F}% secondary panel tint
\definecolor{qaline}{HTML}{2B303D}    % hairline rules / table rules
\definecolor{qatext}{HTML}{ECEEF2}    % primary text (off-white)
\definecolor{qamuted}{HTML}{8B93A6}   % secondary / caption text
\definecolor{qaaccent}{HTML}{D93ED7}  % primary accent -- violet du logo IUT, eclairci pour fond sombre
\definecolor{qaaccent2}{HTML}{F2B84B} % secondary accent -- or/ambre
\definecolor{qacodekw}{HTML}{C792EA}  % code : mots-cles
\definecolor{qacodestr}{HTML}{C3E88D} % code : chaines
\definecolor{qacodecom}{HTML}{7F8AA3} % code : commentaires
\definecolor{qacodeann}{HTML}{F2B84B} % code : annotations / javadoc

\setbeamercolor{background canvas}{bg=qabg}
\setbeamercolor{normal text}{fg=qatext}
\setbeamercolor{palette primary}{bg=qabg, fg=qamuted}
\setbeamercolor{frametitle}{bg=qabg, fg=qatext}
\setbeamercolor{progress bar}{fg=qaaccent, bg=qaline}
\setbeamercolor{alerted text}{fg=qaaccent2}
\setbeamercolor{block title}{bg=qapanel, fg=qaaccent}
\setbeamercolor{block body}{bg=qapanel, fg=qatext}
\setbeamercolor{block title alerted}{bg=qapanel, fg=qaaccent2}
\setbeamercolor{block body alerted}{bg=qapanel, fg=qatext}
\setbeamercolor{block title example}{bg=qapanelalt, fg=qaaccent}
\setbeamercolor{block body example}{bg=qapanelalt, fg=qatext}
\setbeamercolor{example text}{fg=qaaccent}
\setbeamercolor{itemize item}{fg=qaaccent}
\setbeamercolor{itemize subitem}{fg=qaaccent2}
\setbeamercolor{enumerate item}{fg=qaaccent}
\setbeamercolor{enumerate subitem}{fg=qaaccent2}
\setbeamercolor{caption}{fg=qamuted}
\setbeamercolor{page number in head/foot}{fg=qamuted, bg=qabg}
\setbeamercolor{title}{fg=qatext}
\setbeamercolor{subtitle}{fg=qamuted}
\setbeamercolor{author}{fg=qatext}
\setbeamercolor{institute}{fg=qamuted}
\setbeamercolor{date}{fg=qamuted}
\setbeamercolor{section in toc}{fg=qatext}
\setbeamercolor{subsection in toc}{fg=qamuted}

\setbeamertemplate{blocks}[rounded][shadow=false]
\setbeamertemplate{itemize item}{\raisebox{0.15em}{\scriptsize\textcolor{qaaccent}{\ensuremath{\blacktriangleright}}}}
\setbeamertemplate{itemize subitem}{\raisebox{0.1em}{\tiny\textcolor{qaaccent2}{\ensuremath{\bullet}}}}

\setbeamerfont{frametitle}{series=\bfseries, size=\Large}
\setbeamerfont{title}{series=\bfseries, size=\Large}
\setbeamerfont{subtitle}{size=\normalsize}
\setbeamerfont{institute}{size=\scriptsize}
\setbeamerfont{date}{size=\scriptsize}
\setbeamerfont{author}{size=\small}

\metroset{sectionpage=none, subsectionpage=none, numbering=fraction, progressbar=frametitle}
\setbeamertemplate{section in toc}[sections numbered]

% ---------------------------------------------------------------
% LISTINGS -- theme sombre (lstlisting reste utilisable tel quel)
% \lstcode : police + couleur de base, a reutiliser dans les
%            surcharges locales (ex. basicstyle=\lstcode\scriptsize)
% ---------------------------------------------------------------
\newcommand{\lstcode}{\ttfamily\color{qatext}}

\lstdefinelanguage{TypeScript}{
  keywords={typeof, new, true, false, catch, function, return, null, switch, var,
    if, in, while, do, else, case, break, const, let, interface, class, extends,
    implements, enum, export, import, this, super, constructor, public, private,
    protected, static, async, await, throw, throws, try, finally, void},
  ndkeywords={boolean, any, number, string, int, long, double},
  sensitive=true,
  comment=[l]{//},
  morecomment=[s]{/*}{*/},
  morestring=[b]',
  morestring=[b]"
}

\lstset{
  language=Java,
  basicstyle=\lstcode\footnotesize,
  keywordstyle=\color{qacodekw}\bfseries,
  ndkeywordstyle=\color{qaaccent2}\bfseries,
  identifierstyle=\color{qatext},
  stringstyle=\color{qacodestr},
  commentstyle=\color{qacodecom}\itshape,
  morecomment=[s][\color{qacodeann}]{/**}{*/},
  backgroundcolor=\color{qapanelalt},
  rulecolor=\color{qaline},
  frame=single,
  framesep=3pt,
  xleftmargin=4pt,
  xrightmargin=4pt,
  keepspaces=true,
  columns=fullflexible,
  showstringspaces=false,
  breaklines=true,
  tabsize=2,
  escapeinside={(*@}{@*)},
}

% ---------------------------------------------------------------
% CARD-STYLE BOXES
% ---------------------------------------------------------------
\tcbset{
  figcard/.style={colback=white, colframe=white, boxrule=0pt, arc=4pt,
    left=3pt, right=3pt, top=3pt, bottom=3pt, boxsep=0pt,
    halign=center, valign=center},
  datacard/.style={colback=qapanel, colframe=qaline, boxrule=0.4pt, arc=3pt,
    left=8pt, right=8pt, top=6pt, bottom=6pt, coltext=qatext,
    before upper={\color{qatext}}, halign=center},
  codebox/.style={colback=qapanelalt, colframe=qaline, boxrule=0.4pt, arc=3pt,
    left=8pt, right=8pt, top=5pt, bottom=5pt, listing only,
    listing options={basicstyle=\ttfamily\small\color{qaaccent},
      breaklines=true, showstringspaces=false, columns=fullflexible,
      keepspaces=true, frame=none, backgroundcolor={}}}
}
\newtcolorbox{figcard}[1][]{figcard,#1}
\newtcolorbox{datacard}{datacard}
\newtcblisting{codebox}{codebox}

% ---------------------------------------------------------------
% HANDY MACROS
% ---------------------------------------------------------------
\newcommand{\hl}[1]{\alert{#1}}
\newcommand{\source}[1]{{\scriptsize\color{qamuted}#1}}

% Sommaire automatique en debut de section
\AtBeginSection[]{%
  \begin{frame}<beamer>{Plan du cours}
    \tableofcontents[currentsection]
  \end{frame}}

% Page de titre : \rtitleslide{numero de chapitre}{titre du chapitre}
% Les auteurs sont repris de \rauthors (modifiable avant \begin{document})
\newcommand{\rauthors}{Selena Lamari et Chouki Tibermacine}
\newcommand{\rtitleslide}[2]{%
\begin{frame}[plain,noframenumbering]
  \begin{center}
    \vspace{2mm}
    {\color{qaaccent}\rule{2.1cm}{1.2pt}}\\[6mm]
    {\Large\bfseries\color{qatext} R3.07}\\[1mm]
    {\Large\bfseries\color{qatext} SQL dans un langage de programmation}\\[5mm]
    {\large\color{qamuted} Chapitre #1}\\[1mm]
    {\large\color{qaaccent} #2}\\[6mm]
    \includegraphics[height=2.1cm]{iut_logo_transparent.png}

    \vspace{4mm}

    {\small\color{qatext} \rauthors}\\[1mm]
    {\scriptsize\color{qamuted} BUT 2 Informatique --- IUT de Vannes --- Universit\'e Bretagne Sud}
  \end{center}
\end{frame}
}

% ---- Diapo de fin « Questions ? » (même modèle que R5.A.08) ----
\newcommand{\rcontact}{Chouki.Tibermacine@univ-ubs.fr}
\newcommand{\rquestionslide}{%
\begin{frame}{Questions ?}
  \begin{center}
    \vspace{2mm}
    \begin{tikzpicture}
      \node[regular polygon, regular polygon sides=6, draw=qaline, thick,
        minimum size=4.6cm, rotate=90] (hexin) {};
      \node[regular polygon, regular polygon sides=6, draw=qaaccent, dashed,
        thick, minimum size=6.2cm, rotate=90] (hexout) {};
      \node[font=\fontsize{92}{92}\selectfont\bfseries, color=qaaccent2] at (hexin) {?};
    \end{tikzpicture}\\[5mm]
    {\small\color{qamuted} \rcontact}
  \end{center}
\end{frame}}

%   \graphicspath{{../../latex/common/assets/}{img/}{figs/}}
% ---------------------------------------------------------------
\usetheme{metropolis}
\usepackage{appendixnumberbeamer}

\usepackage[T1]{fontenc}
\usepackage[utf8]{inputenc}
\usepackage[french]{babel}
\usepackage[scale=0.94]{tgheros}
\renewcommand{\familydefault}{\sfdefault}
\usepackage{amsmath,amssymb,amsfonts}
\usepackage{graphicx}
\usepackage{booktabs}
\usepackage{multirow}
\usepackage{tikz}
\usetikzlibrary{shapes.geometric, shapes, arrows, arrows.meta, positioning, calc, fit}
\usepackage{xcolor}
\usepackage{listings}
\usepackage[skins,listings]{tcolorbox}
\usepackage{qrcode}
\usepackage{hyperref}

% ---------------------------------------------------------------
% PALETTE (dark-mode tech: near-black canvas, violet (logo IUT) + gold accents)
% ---------------------------------------------------------------
\definecolor{qabg}{HTML}{0E1016}      % canvas background
\definecolor{qapanel}{HTML}{181B24}   % card / block panel
\definecolor{qapanelalt}{HTML}{20242F}% secondary panel tint
\definecolor{qaline}{HTML}{2B303D}    % hairline rules / table rules
\definecolor{qatext}{HTML}{ECEEF2}    % primary text (off-white)
\definecolor{qamuted}{HTML}{8B93A6}   % secondary / caption text
\definecolor{qaaccent}{HTML}{D93ED7}  % primary accent -- violet du logo IUT, eclairci pour fond sombre
\definecolor{qaaccent2}{HTML}{F2B84B} % secondary accent -- or/ambre
\definecolor{qacodekw}{HTML}{C792EA}  % code : mots-cles
\definecolor{qacodestr}{HTML}{C3E88D} % code : chaines
\definecolor{qacodecom}{HTML}{7F8AA3} % code : commentaires
\definecolor{qacodeann}{HTML}{F2B84B} % code : annotations / javadoc

\setbeamercolor{background canvas}{bg=qabg}
\setbeamercolor{normal text}{fg=qatext}
\setbeamercolor{palette primary}{bg=qabg, fg=qamuted}
\setbeamercolor{frametitle}{bg=qabg, fg=qatext}
\setbeamercolor{progress bar}{fg=qaaccent, bg=qaline}
\setbeamercolor{alerted text}{fg=qaaccent2}
\setbeamercolor{block title}{bg=qapanel, fg=qaaccent}
\setbeamercolor{block body}{bg=qapanel, fg=qatext}
\setbeamercolor{block title alerted}{bg=qapanel, fg=qaaccent2}
\setbeamercolor{block body alerted}{bg=qapanel, fg=qatext}
\setbeamercolor{block title example}{bg=qapanelalt, fg=qaaccent}
\setbeamercolor{block body example}{bg=qapanelalt, fg=qatext}
\setbeamercolor{example text}{fg=qaaccent}
\setbeamercolor{itemize item}{fg=qaaccent}
\setbeamercolor{itemize subitem}{fg=qaaccent2}
\setbeamercolor{enumerate item}{fg=qaaccent}
\setbeamercolor{enumerate subitem}{fg=qaaccent2}
\setbeamercolor{caption}{fg=qamuted}
\setbeamercolor{page number in head/foot}{fg=qamuted, bg=qabg}
\setbeamercolor{title}{fg=qatext}
\setbeamercolor{subtitle}{fg=qamuted}
\setbeamercolor{author}{fg=qatext}
\setbeamercolor{institute}{fg=qamuted}
\setbeamercolor{date}{fg=qamuted}
\setbeamercolor{section in toc}{fg=qatext}
\setbeamercolor{subsection in toc}{fg=qamuted}

\setbeamertemplate{blocks}[rounded][shadow=false]
\setbeamertemplate{itemize item}{\raisebox{0.15em}{\scriptsize\textcolor{qaaccent}{\ensuremath{\blacktriangleright}}}}
\setbeamertemplate{itemize subitem}{\raisebox{0.1em}{\tiny\textcolor{qaaccent2}{\ensuremath{\bullet}}}}

\setbeamerfont{frametitle}{series=\bfseries, size=\Large}
\setbeamerfont{title}{series=\bfseries, size=\Large}
\setbeamerfont{subtitle}{size=\normalsize}
\setbeamerfont{institute}{size=\scriptsize}
\setbeamerfont{date}{size=\scriptsize}
\setbeamerfont{author}{size=\small}

\metroset{sectionpage=none, subsectionpage=none, numbering=fraction, progressbar=frametitle}
\setbeamertemplate{section in toc}[sections numbered]

% ---------------------------------------------------------------
% LISTINGS -- theme sombre (lstlisting reste utilisable tel quel)
% \lstcode : police + couleur de base, a reutiliser dans les
%            surcharges locales (ex. basicstyle=\lstcode\scriptsize)
% ---------------------------------------------------------------
\newcommand{\lstcode}{\ttfamily\color{qatext}}

\lstdefinelanguage{TypeScript}{
  keywords={typeof, new, true, false, catch, function, return, null, switch, var,
    if, in, while, do, else, case, break, const, let, interface, class, extends,
    implements, enum, export, import, this, super, constructor, public, private,
    protected, static, async, await, throw, throws, try, finally, void},
  ndkeywords={boolean, any, number, string, int, long, double},
  sensitive=true,
  comment=[l]{//},
  morecomment=[s]{/*}{*/},
  morestring=[b]',
  morestring=[b]"
}

\lstset{
  language=Java,
  basicstyle=\lstcode\footnotesize,
  keywordstyle=\color{qacodekw}\bfseries,
  ndkeywordstyle=\color{qaaccent2}\bfseries,
  identifierstyle=\color{qatext},
  stringstyle=\color{qacodestr},
  commentstyle=\color{qacodecom}\itshape,
  morecomment=[s][\color{qacodeann}]{/**}{*/},
  backgroundcolor=\color{qapanelalt},
  rulecolor=\color{qaline},
  frame=single,
  framesep=3pt,
  xleftmargin=4pt,
  xrightmargin=4pt,
  keepspaces=true,
  columns=fullflexible,
  showstringspaces=false,
  breaklines=true,
  tabsize=2,
  escapeinside={(*@}{@*)},
}

% ---------------------------------------------------------------
% CARD-STYLE BOXES
% ---------------------------------------------------------------
\tcbset{
  figcard/.style={colback=white, colframe=white, boxrule=0pt, arc=4pt,
    left=3pt, right=3pt, top=3pt, bottom=3pt, boxsep=0pt,
    halign=center, valign=center},
  datacard/.style={colback=qapanel, colframe=qaline, boxrule=0.4pt, arc=3pt,
    left=8pt, right=8pt, top=6pt, bottom=6pt, coltext=qatext,
    before upper={\color{qatext}}, halign=center},
  codebox/.style={colback=qapanelalt, colframe=qaline, boxrule=0.4pt, arc=3pt,
    left=8pt, right=8pt, top=5pt, bottom=5pt, listing only,
    listing options={basicstyle=\ttfamily\small\color{qaaccent},
      breaklines=true, showstringspaces=false, columns=fullflexible,
      keepspaces=true, frame=none, backgroundcolor={}}}
}
\newtcolorbox{figcard}[1][]{figcard,#1}
\newtcolorbox{datacard}{datacard}
\newtcblisting{codebox}{codebox}

% ---------------------------------------------------------------
% HANDY MACROS
% ---------------------------------------------------------------
\newcommand{\hl}[1]{\alert{#1}}
\newcommand{\source}[1]{{\scriptsize\color{qamuted}#1}}

% Sommaire automatique en debut de section
\AtBeginSection[]{%
  \begin{frame}<beamer>{Plan du cours}
    \tableofcontents[currentsection]
  \end{frame}}

% Page de titre : \rtitleslide{numero de chapitre}{titre du chapitre}
% Les auteurs sont repris de \rauthors (modifiable avant \begin{document})
\newcommand{\rauthors}{Selena Lamari et Chouki Tibermacine}
\newcommand{\rtitleslide}[2]{%
\begin{frame}[plain,noframenumbering]
  \begin{center}
    \vspace{2mm}
    {\color{qaaccent}\rule{2.1cm}{1.2pt}}\\[6mm]
    {\Large\bfseries\color{qatext} R3.07}\\[1mm]
    {\Large\bfseries\color{qatext} SQL dans un langage de programmation}\\[5mm]
    {\large\color{qamuted} Chapitre #1}\\[1mm]
    {\large\color{qaaccent} #2}\\[6mm]
    \includegraphics[height=2.1cm]{iut_logo_transparent.png}

    \vspace{4mm}

    {\small\color{qatext} \rauthors}\\[1mm]
    {\scriptsize\color{qamuted} BUT 2 Informatique --- IUT de Vannes --- Universit\'e Bretagne Sud}
  \end{center}
\end{frame}
}

% ---- Diapo de fin « Questions ? » (même modèle que R5.A.08) ----
\newcommand{\rcontact}{Chouki.Tibermacine@univ-ubs.fr}
\newcommand{\rquestionslide}{%
\begin{frame}{Questions ?}
  \begin{center}
    \vspace{2mm}
    \begin{tikzpicture}
      \node[regular polygon, regular polygon sides=6, draw=qaline, thick,
        minimum size=4.6cm, rotate=90] (hexin) {};
      \node[regular polygon, regular polygon sides=6, draw=qaaccent, dashed,
        thick, minimum size=6.2cm, rotate=90] (hexout) {};
      \node[font=\fontsize{92}{92}\selectfont\bfseries, color=qaaccent2] at (hexin) {?};
    \end{tikzpicture}\\[5mm]
    {\small\color{qamuted} \rcontact}
  \end{center}
\end{frame}}

\graphicspath{{../../latex/common/assets/}}

\title{Programmer dans la base : proc\'edures, exceptions, curseurs, triggers}
\author[C. Tibermacine]{Selena Lamari et Chouki Tibermacine}
\date{}

% Sortie de psql
\lstdefinestyle{out}{language={}, basicstyle=\lstcode\scriptsize, backgroundcolor=\color{qapanel}}

\begin{document}

\rtitleslide{7}{Programmer dans la base : proc\'edures, exceptions, curseurs, triggers}

\begin{frame}{Avant de commencer}
  \begin{alertblock}{La syntaxe de base est dans le TP6}
    Le \textbf{mini-tutoriel PL/pgSQL} du sujet de TP6 pr\'esente la syntaxe : structure d'une proc\'edure,
    \texttt{DECLARE}, \texttt{SELECT \ldots{} INTO}, \texttt{IF}, \texttt{RAISE EXCEPTION}, curseur simple,
    fonction de trigger, \texttt{OLD}/\texttt{NEW}, \texttt{CREATE TRIGGER}. Le sujet liste aussi des
    exemples de proc\'edures et de triggers pour l'application bancaire.
  \end{alertblock}
  \vspace{2mm}
  Ce chapitre ne la r\'ep\`ete pas. Il r\'epond aux questions que le tutoriel laisse ouvertes :
  \begin{itemize}
    \item \textbf{faut-il} mettre cette logique dans la base ?
    \item fonction ou proc\'edure ? peut-on valider (\texttt{COMMIT}) dans une proc\'edure ?
    \item comment \textbf{rattraper} une erreur en PL/pgSQL, et comment Java la re\c{c}oit-il ?
    \item quand un curseur est-il utile, c\^ot\'e base et c\^ot\'e JDBC ?
    \item quels triggers au-del\`a du \texttt{FOR EACH ROW} simple, et quels pi\`eges ?
    \item comment s\'ecuriser et versionner ce code ?
  \end{itemize}
\end{frame}

% ================================================================
\section{O\`u placer la logique ?}

\begin{frame}{Java ou base de donn\'ees : les crit\`eres}
  \begin{datacard}
    \small
    \begin{tabular}{@{}p{3.6cm}p{4.3cm}p{4.3cm}@{}}
      \toprule
      \textbf{Crit\`ere} & \textbf{Plut\^ot dans la base} & \textbf{Plut\^ot en Java} \\
      \midrule
      Clients de la base & plusieurs applications, scripts, acc\`es SQL direct & une seule application \\
      Volume trait\'e & beaucoup de lignes, peu de r\'esultat & peu de lignes \\
      Allers-retours & nombreux (boucle de requ\^etes) & un ou deux \\
      Tests, d\'ebogage & --- & outils Java (JUnit, IDE) \\
      Mont\'ee en charge & --- & on ajoute des serveurs d'application ; la base est plus dure \`a multiplier \\
      Portabilit\'e & --- & PL/pgSQL n'existe que sur PostgreSQL \\
      \bottomrule
    \end{tabular}
  \end{datacard}
  \vspace{1mm}
  {\small\color{qamuted} Le sujet du TP6 liste les avantages des proc\'edures et triggers ; ce tableau rappelle qu'ils ont aussi un prix.}
\end{frame}

\begin{frame}[fragile]{La r\`egle de l'outil le plus simple}
  \begin{enumerate}
    \item \textbf{Contrainte d\'eclarative} d'abord : \texttt{NOT NULL}, \texttt{CHECK}, \texttt{UNIQUE}, cl\'e \'etrang\`ere
    \item \textbf{Trigger} si la r\`egle d\'epend d'autres lignes ou d'autres tables, ou s'il faut une action automatique (audit)
    \item \textbf{Proc\'edure} pour une op\'eration m\'etier compl\`ete et nomm\'ee (virement, cl\^oture)
  \end{enumerate}
  \vspace{2mm}
  \begin{exampleblock}{Exemple : \og interdire un solde n\'egatif \fg{} (liste d'exemples du TP6)}
    Un trigger fonctionne, mais une contrainte suffit, est plus rapide et ne peut pas \^etre oubli\'ee :
\begin{lstlisting}[language=SQL]
ALTER TABLE Compte ADD CONSTRAINT chk_solde CHECK (solde >= 0);
\end{lstlisting}
    Sa violation renvoie un code pr\'ecis (\texttt{23514}) et le nom de la contrainte \`a Java.
  \end{exampleblock}
\end{frame}

% ================================================================
\section{Fonctions et proc\'edures}

\begin{frame}{Fonction ou proc\'edure ?}
  \begin{datacard}
    \small
    \begin{tabular}{@{}p{3.3cm}p{4.6cm}p{4.6cm}@{}}
      \toprule
       & \textbf{\texttt{FUNCTION}} & \textbf{\texttt{PROCEDURE}} (PostgreSQL 11+) \\
      \midrule
      Appel & dans une requ\^ete : \texttt{SELECT f(\ldots)}, \texttt{FROM f(\ldots)}, \texttt{WHERE} & \texttt{CALL p(\ldots)} uniquement \\
      R\'esultat & une valeur, une ligne ou une table (\texttt{RETURNS TABLE}, \texttt{SETOF}) & aucun, ou param\`etres \texttt{OUT}/\texttt{INOUT} \\
      \texttt{COMMIT}/\texttt{ROLLBACK} internes & \alert{interdits} & autoris\'es (sous conditions) \\
      Usage typique & calcul, requ\^ete param\'etr\'ee, fonction de trigger & op\'eration m\'etier, traitement par lots \\
      \bottomrule
    \end{tabular}
  \end{datacard}
  \vspace{2mm}
  \begin{itemize}
    \item Le relev\'e de compte du TP6 est une \textbf{fonction} (on l'interroge avec \texttt{SELECT}) ; le virement est une \textbf{proc\'edure}
  \end{itemize}
\end{frame}

\begin{frame}[fragile]{Valider par lots dans une proc\'edure}
  \begin{columns}[T]
    \begin{column}{0.55\textwidth}
\begin{lstlisting}[language=SQL,basicstyle=\lstcode\scriptsize]
CREATE PROCEDURE verser_interets(
    p_taux NUMERIC, p_lot INT DEFAULT 1000)
LANGUAGE plpgsql AS $$
DECLARE c RECORD; n INT := 0;
BEGIN
  FOR c IN SELECT numCompte FROM Compte
            WHERE typeCompte = 'EPARGNE' LOOP
    UPDATE Compte
       SET solde = round(solde * (1 + p_taux), 2)
     WHERE numCompte = c.numCompte;
    n := n + 1;
    IF n % p_lot = 0 THEN
      COMMIT;                   -- lot valide
    END IF;
  END LOOP;
END $$;
\end{lstlisting}
    \end{column}
    \begin{column}{0.42\textwidth}
      \begin{itemize}
        \item Traitement \textit{batch} (le PN cite les \og batch et proc\'edures stock\'ees \fg{})
        \item Un incident en cours de route ne fait perdre que le lot courant
        \item Les verrous sont rel\^ach\'es \`a chaque \texttt{COMMIT}
      \end{itemize}
      \vspace{1mm}
      \textbf{Interdit} (test\'e) :
      \begin{itemize}
        \item \texttt{COMMIT} dans une fonction
        \item \texttt{CALL} dans un \texttt{BEGIN} explicite
      \end{itemize}
\begin{lstlisting}[style=out]
ERROR: invalid transaction termination
\end{lstlisting}
    \end{column}
  \end{columns}
\end{frame}

\begin{frame}[fragile]{Cons\'equence c\^ot\'e Java}
  \begin{itemize}
    \item En \texttt{autoCommit = false}, le pilote JDBC ouvre lui-m\^eme une transaction avant chaque requ\^ete
    \item Une proc\'edure qui fait \texttt{COMMIT} doit donc \^etre appel\'ee en \textbf{autocommit}
  \end{itemize}
\begin{lstlisting}[language=Java]
c.setAutoCommit(true);
try (Statement s = c.createStatement()) {
  s.execute("CALL verser_interets(0.02)");                // OK
}
c.setAutoCommit(false);
try (Statement s = c.createStatement()) {
  s.execute("CALL verser_interets(0.02)");                // SQLState 2D000
}                                                         // invalid transaction termination
\end{lstlisting}
  \begin{itemize}
    \item Une proc\'edure \textbf{sans} \texttt{COMMIT} interne (le virement) s'appelle normalement dans la transaction de l'application
  \end{itemize}
\end{frame}

\begin{frame}[fragile]{Volatilit\'e d'une fonction}
  \begin{itemize}
    \item \texttt{VOLATILE} (d\'efaut) : r\'esultat impr\'evisible, r\'e\'evalu\'ee \`a chaque ligne (\texttt{random()}, \texttt{now()} entre deux transactions, tout ce qui \'ecrit)
    \item \texttt{STABLE} : m\^eme r\'esultat pendant une requ\^ete, peut lire la base
    \item \texttt{IMMUTABLE} : ne d\'epend que de ses arguments
  \end{itemize}
  Lien avec le chapitre 3 : un \textbf{index sur expression} exige une fonction \texttt{IMMUTABLE}
\begin{lstlisting}[language=SQL]
CREATE FUNCTION categorie(s NUMERIC) RETURNS TEXT LANGUAGE plpgsql
  AS $$ BEGIN RETURN CASE WHEN s > 1000 THEN 'haut' ELSE 'bas' END; END $$;
CREATE INDEX ON Compte (categorie(solde));
\end{lstlisting}
\begin{lstlisting}[style=out]
ERROR:  functions in index expression must be marked IMMUTABLE
\end{lstlisting}
  {\small D\'eclarer la bonne volatilit\'e aide aussi l'optimiseur. Mentir (\texttt{IMMUTABLE} sur une fonction qui lit une table) donne des index faux.}
\end{frame}

\begin{frame}[fragile]{Appeler fonctions et proc\'edures depuis Java}
\begin{lstlisting}[language=Java]
// Fonction qui renvoie une table : une requete comme une autre
PreparedStatement ps = c.prepareStatement("SELECT * FROM releve_compte(?)");

// Procedure : CALL, en convertissant les types si besoin
PreparedStatement ps = c.prepareStatement("CALL virement(?, ?, CAST(? AS REAL))");
\end{lstlisting}
  \begin{itemize}
    \item PostgreSQL choisit la proc\'edure d'apr\`es les \textbf{types} des arguments ; \texttt{setDouble} envoie un
      \texttt{double precision} : sans \texttt{CAST}, \texttt{procedure virement(integer, integer, double precision) does not exist}
    \item Avec Spring Data : \texttt{@Query(value = "CALL \ldots", nativeQuery = true)} (chapitre 6)
    \item \texttt{@Procedure} (Spring Data) passe les param\`etres \textbf{par nom} : ils doivent porter le m\^eme nom
      en Java et en SQL, sinon l'appel \'echoue (test\'e)
  \end{itemize}
\end{frame}

% ================================================================
\section{Exceptions}

\begin{frame}[fragile]{Lever une erreur exploitable}
\begin{lstlisting}[language=SQL]
RAISE EXCEPTION 'Solde insuffisant sur le compte %', p_num
  USING ERRCODE = 'P0002',
        DETAIL  = format('solde = %s, demande = %s', v_solde, p_montant),
        HINT    = 'Reduire le montant';
\end{lstlisting}
\begin{lstlisting}[style=out]
ERROR:  Solde insuffisant sur le compte 3
DETAIL:  solde = 10, demande = 50
HINT:  Reduire le montant
\end{lstlisting}
  \begin{itemize}
    \item Niveaux : \texttt{DEBUG}, \texttt{LOG}, \texttt{INFO}, \texttt{NOTICE}, \texttt{WARNING} (simple message) ; \texttt{EXCEPTION} (arr\^ete et annule)
    \item \texttt{ERRCODE} : un code par erreur m\'etier, pour que Java les distingue sans analyser le texte.
      Par d\'efaut : \texttt{P0001}
  \end{itemize}
\end{frame}

\begin{frame}[fragile]{Rattraper une erreur : le bloc \texttt{EXCEPTION}}
  \begin{columns}[T]
    \begin{column}{0.55\textwidth}
\begin{lstlisting}[language=SQL,basicstyle=\lstcode\scriptsize]
BEGIN
  INSERT INTO journal(msg) VALUES ('avant');
  BEGIN                                   -- sous-bloc
    INSERT INTO journal(msg) VALUES ('dans le bloc');
    UPDATE Compte SET solde = -5 WHERE numCompte = 1;
  EXCEPTION
    WHEN check_violation THEN
      INSERT INTO journal(msg)
      VALUES ('rattrape : ' || SQLSTATE);
  END;
END;
\end{lstlisting}
\begin{lstlisting}[style=out]
 avant
 rattrape : 23514
\end{lstlisting}
    \end{column}
    \begin{column}{0.42\textwidth}
      \begin{itemize}
        \item Conditions : nom (\texttt{unique\_violation}, \texttt{check\_violation}\ldots), \texttt{SQLSTATE '23505'}, ou \texttt{OTHERS}
        \item \alert{Tout ce que le sous-bloc a fait est annul\'e} : la ligne \og dans le bloc \fg{} a disparu
        \item Le bloc pose un point de sauvegarde : il co\^ute cher, \`a \'eviter dans une boucle sur des milliers de lignes
        \item Ne pas tout avaler avec \texttt{OTHERS} : journaliser puis relancer (\texttt{RAISE;})
      \end{itemize}
    \end{column}
  \end{columns}
\end{frame}

\begin{frame}[fragile]{Inspecter l'erreur rattrap\'ee}
\begin{lstlisting}[language=SQL]
EXCEPTION WHEN OTHERS THEN
  GET STACKED DIAGNOSTICS v_etat       = RETURNED_SQLSTATE,
                          v_msg        = MESSAGE_TEXT,
                          v_contrainte = CONSTRAINT_NAME;
  RAISE NOTICE 'etat=% contrainte=% msg=%', v_etat, v_contrainte, v_msg;
  RAISE;                                       -- relance la meme erreur
\end{lstlisting}
\begin{lstlisting}[style=out]
NOTICE:  etat=23505 contrainte=compte_pkey msg=duplicate key value violates unique constraint "compte_pkey"
\end{lstlisting}
  \begin{itemize}
    \item Utile pour journaliser dans une table d'erreurs avant de relancer
    \item Autres champs : \texttt{PG\_EXCEPTION\_DETAIL}, \texttt{PG\_EXCEPTION\_HINT}, \texttt{TABLE\_NAME}, \texttt{COLUMN\_NAME}\ldots
  \end{itemize}
\end{frame}

\begin{frame}[fragile]{L'erreur vue de Java}
\begin{lstlisting}[language=Java]
} catch (PSQLException e) {                       // org.postgresql.util
  ServerErrorMessage m = e.getServerErrorMessage();
  switch (e.getSQLState()) {
    case "P0002" -> throw new SoldeInsuffisantException(m.getDetail());
    case "23505" -> throw new DoublonException(m.getConstraint());
    case "23514" -> throw new RegleMetierException(m.getConstraint());
    default      -> throw new DataAccessException(e);
  }
}
\end{lstlisting}
  \begin{datacard}
    \scriptsize
    \begin{tabular}{@{}llll@{}}
      \toprule
      \texttt{23505} unicit\'e & \texttt{23503} cl\'e \'etrang\`ere & \texttt{23514} \texttt{CHECK} & \texttt{23502} \texttt{NOT NULL} \\
      \texttt{40001} s\'erialisation (ch.\,2) & \texttt{40P01} interblocage (ch.\,2) & \texttt{P0001} \texttt{RAISE} par d\'efaut & \texttt{P0002}\ldots{} vos codes \\
      \bottomrule
    \end{tabular}
  \end{datacard}
  {\small Test\'e : \texttt{getDetail()} renvoie le \texttt{DETAIL} du \texttt{RAISE} (\og solde = 10 \fg), \texttt{getConstraint()} le nom de la contrainte viol\'ee (\og chk\_solde \fg).}
\end{frame}

% ================================================================
\section{Curseurs}

\begin{frame}{Curseurs c\^ot\'e base : quand ?}
  \begin{itemize}
    \item Un curseur traite le r\'esultat \textbf{ligne par ligne} sans le charger en entier
    \item \texttt{FOR ligne IN SELECT \ldots{} LOOP} ouvre un curseur implicite : c'est la forme \`a privil\'egier
    \item Curseur explicite (\texttt{OPEN}/\texttt{FETCH}/\texttt{CLOSE}, vu dans le tutoriel) quand on veut contr\^oler le parcours :
      plusieurs curseurs en parall\`ele, arr\^et anticip\'e, curseur param\'etr\'e (\texttt{CURSOR (p INT) FOR \ldots})
  \end{itemize}
  \begin{block}{Mais d'abord : peut-on l'\'ecrire en une requ\^ete ?}
    \begin{itemize}
      \item Exercice 3 du TP6 : le relev\'e avec curseur et la version \texttt{sum(\ldots) OVER (\ldots)} donnent le m\^eme r\'esultat
      \item Une requ\^ete ensembliste laisse l'optimiseur choisir (index, parall\'elisme) ; une boucle impose l'ordre et le d\'etail
      \item Curseur justifi\'e : d\'ecision diff\'erente pour chaque ligne, appel d'autres proc\'edures, validation par lots
    \end{itemize}
  \end{block}
\end{frame}

\begin{frame}[fragile]{Curseurs c\^ot\'e JDBC : lire un gros r\'esultat}
  \begin{itemize}
    \item Par d\'efaut, le pilote PostgreSQL \textbf{charge tout le r\'esultat en m\'emoire} avant le premier \texttt{next()}
    \item Avec \texttt{setFetchSize(n)}, il utilise un curseur serveur et lit par paquets de $n$ lignes\ldots{}
      \alert{seulement si l'autocommit est d\'esactiv\'e}
  \end{itemize}
\begin{lstlisting}[language=Java]
c.setAutoCommit(false);                     // indispensable
try (PreparedStatement ps = c.prepareStatement("SELECT id, txt FROM big")) {
  ps.setFetchSize(1000);
  try (ResultSet rs = ps.executeQuery()) { while (rs.next()) { ... } }
}
c.commit();
\end{lstlisting}
  \begin{datacard}
    \small
    \begin{tabular}{@{}lrr@{}}
      \toprule
      \textbf{2 millions de lignes, \texttt{fetchSize = 1000}} & \textbf{1\textsuperscript{re} ligne re\c{c}ue} & \textbf{m\'emoire} \\
      \midrule
      \texttt{autoCommit = true} (fetchSize ignor\'e) & 458 ms & tout le r\'esultat \\
      \texttt{autoCommit = false} & \alert{1 ms} & un paquet \\
      \bottomrule
    \end{tabular}
  \end{datacard}
\end{frame}

% ================================================================
\section{Triggers : au-del\`a des bases}

\begin{frame}{Les variantes de triggers}
  \begin{datacard}
    \small
    \begin{tabular}{@{}p{2.6cm}p{5cm}p{5cm}@{}}
      \toprule
       & \textbf{\texttt{BEFORE}} & \textbf{\texttt{AFTER}} \\
      \midrule
      \texttt{FOR EACH ROW} & peut \textbf{modifier} \texttt{NEW} (valeur calcul\'ee) ou renvoyer \texttt{NULL} (ligne ignor\'ee) & voit la ligne finale ; valeur de retour ignor\'ee (audit, propagation) \\
      \texttt{FOR EACH STATEMENT} & une fois par instruction, avant & une fois par instruction ; acc\`es \`a toutes les lignes via les \textbf{tables de transition} \\
      \bottomrule
    \end{tabular}
  \end{datacard}
  \vspace{2mm}
  \begin{itemize}
    \item Filtres : \texttt{UPDATE OF colonne} (seulement si cette colonne est cit\'ee), \texttt{WHEN (condition)} (\'evalu\'ee avant d'appeler la fonction)
    \item Plusieurs triggers sur le m\^eme \'ev\'enement : ex\'ecut\'es par \textbf{ordre alphab\'etique} de nom
    \item Variables disponibles : \texttt{TG\_OP} (\texttt{INSERT}/\texttt{UPDATE}/\texttt{DELETE}), \texttt{TG\_TABLE\_NAME}, \texttt{TG\_WHEN}\ldots
  \end{itemize}
\end{frame}

\begin{frame}[fragile]{\texttt{BEFORE} : compl\'eter la ligne, avec filtres}
\begin{lstlisting}[language=SQL]
CREATE FUNCTION trg_horodater() RETURNS trigger LANGUAGE plpgsql AS $$
BEGIN
  NEW.maj := now();                  -- on modifie la ligne avant ecriture
  RETURN NEW;
END $$;

CREATE TRIGGER compte_horodater
  BEFORE UPDATE OF solde ON Compte
  FOR EACH ROW
  WHEN (OLD.solde IS DISTINCT FROM NEW.solde)
  EXECUTE FUNCTION trg_horodater();
\end{lstlisting}
  \begin{itemize}
    \item Test\'e : \texttt{SET solde = solde} (inchang\'e) et \texttt{SET type = type} (autre colonne) n'horodatent pas ;
      \texttt{SET solde = solde + 1} horodate
    \item \texttt{WHEN} \'evite m\^eme l'appel de la fonction : moins co\^uteux qu'un \texttt{IF} \`a l'int\'erieur
  \end{itemize}
\end{frame}

\begin{frame}[fragile]{Un trigger d'audit pour toutes les tables}
\begin{lstlisting}[language=SQL,basicstyle=\lstcode\scriptsize]
CREATE TABLE audit(id BIGSERIAL, quand TIMESTAMPTZ DEFAULT now(), utilisateur TEXT DEFAULT current_user,
                   table_ TEXT, op TEXT, avant JSONB, apres JSONB);

CREATE FUNCTION trg_audit() RETURNS trigger LANGUAGE plpgsql AS $$
BEGIN
  INSERT INTO audit(table_, op, avant, apres) VALUES (TG_TABLE_NAME, TG_OP,
    CASE WHEN TG_OP <> 'INSERT' THEN to_jsonb(OLD) END,
    CASE WHEN TG_OP <> 'DELETE' THEN to_jsonb(NEW) END);
  RETURN NULL;
END $$;

CREATE TRIGGER compte_audit AFTER INSERT OR UPDATE OR DELETE ON Compte
  FOR EACH ROW EXECUTE FUNCTION trg_audit();
CREATE TRIGGER client_audit AFTER INSERT OR UPDATE OR DELETE ON Client
  FOR EACH ROW EXECUTE FUNCTION trg_audit();       -- meme fonction
\end{lstlisting}
  {\small Une seule fonction, r\'eutilis\'ee par toutes les tables, contre une fonction par table dans l'exercice 2 du TP6.
  Garde l'\'etat avant/apr\`es, l'auteur et la date (AC24.02 : int\'egrit\'e et tra\c{c}abilit\'e).}
\end{frame}

\begin{frame}[fragile]{Niveau instruction et tables de transition}
\begin{lstlisting}[language=SQL,basicstyle=\lstcode\scriptsize]
CREATE FUNCTION trg_gros_volume() RETURNS trigger LANGUAGE plpgsql AS $$
DECLARE v_nb INT; v_total NUMERIC;
BEGIN
  SELECT count(*), sum(n.solde - o.solde) INTO v_nb, v_total
    FROM nouveaux n JOIN anciens o USING (numCompte);      -- toutes les lignes modifiees
  IF v_nb > 100 THEN
    INSERT INTO alerte(nb, total) VALUES (v_nb, v_total);
  END IF;
  RETURN NULL;
END $$;

CREATE TRIGGER compte_volume AFTER UPDATE ON Compte
  REFERENCING OLD TABLE AS anciens NEW TABLE AS nouveaux
  FOR EACH STATEMENT EXECUTE FUNCTION trg_gros_volume();
\end{lstlisting}
  \begin{itemize}
    \item Un seul appel pour un \texttt{UPDATE} de 10\,000 lignes, au lieu de 10\,000 avec \texttt{FOR EACH ROW}
    \item Permet des contr\^oles globaux : volume, total d\'eplac\'e, coh\'erence entre lignes
  \end{itemize}
\end{frame}

\begin{frame}[fragile]{Les pi\`eges des triggers}
  \begin{itemize}
    \item \textbf{Annulation silencieuse} : un \texttt{BEFORE} qui renvoie \texttt{NULL} ignore la ligne sans erreur
\begin{lstlisting}[style=out]
DELETE FROM Compte WHERE numCompte = 9;   -- DELETE 0, aucune erreur pour l'application
\end{lstlisting}
      Pour refuser une op\'eration : \texttt{RAISE EXCEPTION}, que Java recevra
    \item \textbf{Logique cach\'ee} : le d\'eveloppeur Java ne voit pas qu'un \texttt{UPDATE} \'ecrit dans trois autres tables ; documenter dans le README
    \item \textbf{Cascades} : un trigger qui modifie une table qui a elle-m\^eme des triggers\ldots{}
    \item \textbf{Co\^ut} : un trigger ligne par ligne s'ex\'ecute pour chaque ligne d'un \texttt{UPDATE} de masse
    \item \textbf{JPA} : une colonne modifi\'ee par un trigger n'est pas vue par l'entit\'e en m\'emoire ;
      relire avec \texttt{em.refresh(entite)}, ou d\'eclarer la colonne \texttt{@Generated} (Hibernate)
  \end{itemize}
\end{frame}

% ================================================================
\section{S\'ecuriser et industrialiser}

\begin{frame}[fragile]{N'autoriser que les proc\'edures : \texttt{SECURITY DEFINER}}
\begin{lstlisting}[language=SQL,basicstyle=\lstcode\scriptsize]
CREATE ROLE appli LOGIN PASSWORD '...';          -- compte utilise par l'application Java
GRANT SELECT ON Compte TO appli;                 -- lecture seule sur la table

CREATE PROCEDURE debiter(p_num INT, p_m NUMERIC)
LANGUAGE plpgsql
SECURITY DEFINER                                 -- s'execute avec les droits du proprietaire
SET search_path = public                         -- evite le detournement par un autre schema
AS $$ BEGIN ... END $$;

REVOKE EXECUTE ON PROCEDURE debiter(INT, NUMERIC) FROM PUBLIC;
GRANT  EXECUTE ON PROCEDURE debiter(INT, NUMERIC) TO appli;
\end{lstlisting}
\begin{lstlisting}[style=out]
appli=> UPDATE Compte SET solde = 0 WHERE numCompte = 1;   ERROR: permission denied for table compte
appli=> CALL debiter(1, 10);                               CALL
\end{lstlisting}
  {\small M\^eme si l'application est compromise (injection SQL), elle ne peut modifier les soldes que par les r\`egles de la proc\'edure (AC24.02).}
\end{frame}

\begin{frame}[fragile]{SQL dynamique sans injection}
\begin{lstlisting}[language=SQL]
CREATE FUNCTION compter(p_table TEXT, p_type TEXT) RETURNS BIGINT LANGUAGE plpgsql AS $$
DECLARE n BIGINT;
BEGIN
  EXECUTE format('SELECT count(*) FROM %I WHERE type = $1', p_table)   -- %I : identifiant
     INTO n USING p_type;                                              -- $1 : valeur
  RETURN n;
END $$;
\end{lstlisting}
\begin{lstlisting}[style=out]
SELECT compter('cpt; DROP TABLE cpt', 'x');
ERROR:  relation "cpt; DROP TABLE cpt" does not exist
\end{lstlisting}
  \begin{itemize}
    \item Jamais de concat\'enation \texttt{'... ' || p\_table || ' ...'} : c'est l'injection SQL du chapitre 1, c\^ot\'e serveur
    \item \texttt{\%I} pour les noms de tables et colonnes, \texttt{USING} pour les valeurs (comme un \texttt{PreparedStatement})
  \end{itemize}
\end{frame}

\begin{frame}[fragile]{Versionner le code de la base}
  \begin{itemize}
    \item Le code PL/pgSQL est du code : il va dans le d\'ep\^ot Git, avec des scripts num\'erot\'es
  \end{itemize}
\begin{lstlisting}[language={}]
src/main/resources/db/migration/
  V1__schema.sql
  V2__contraintes.sql
  V3__procedure_virement.sql
  V4__audit.sql
\end{lstlisting}
  \begin{itemize}
    \item Un outil de migration (\textbf{Flyway}, Liquibase) applique au d\'emarrage les scripts pas encore pass\'es, dans l'ordre,
      et garde l'historique dans une table ; avec Spring Boot, il suffit d'ajouter les d\'ependances \texttt{flyway-core} et \texttt{flyway-database-postgresql}
    \item \texttt{CREATE OR REPLACE} pour faire \'evoluer une fonction sans la supprimer
    \item Tester les proc\'edures : tests d'int\'egration Java contre une vraie base PostgreSQL (conteneur Docker)
  \end{itemize}
\end{frame}

\begin{frame}{\`A retenir}
  \begin{itemize}
    \item Contrainte d'abord, trigger ensuite, proc\'edure pour une op\'eration m\'etier nomm\'ee
    \item Fonction = interrogeable en SQL, sans \texttt{COMMIT} ; proc\'edure = \texttt{CALL}, peut valider par lots (en autocommit c\^ot\'e JDBC)
    \item Erreurs : \texttt{RAISE \ldots{} USING ERRCODE} c\^ot\'e base, \texttt{getSQLState()} c\^ot\'e Java ; un bloc \texttt{EXCEPTION} annule son contenu
    \item Curseurs : pr\'ef\'erer une requ\^ete ensembliste ; c\^ot\'e JDBC, \texttt{setFetchSize} + autocommit d\'esactiv\'e
    \item Triggers : \texttt{BEFORE} pour compl\'eter, \texttt{AFTER} pour tracer, niveau instruction pour les gros volumes ; jamais d'annulation silencieuse
    \item S\'ecurit\'e : \texttt{SECURITY DEFINER} + \texttt{GRANT EXECUTE}, \texttt{format('\%I')} + \texttt{USING}
  \end{itemize}
\end{frame}

\rquestionslide

\end{document}
